\section{Dynkin Diagrams}

\begin{frame}{Simple Roots and Dynkin Diagrams}
    To formalize the classification for any rank, we introduce \textbf{Simple Roots} ($\Delta$): a minimal basis of positive roots that generate the entire system.

    \vspace{0.5cm}
    We condense the geometry of simple roots into a \textbf{Dynkin Diagram}:
    \begin{itemize}
        \item \textbf{Nodes:} Represent simple roots.
        \item \textbf{Edges:} Represent the angle between them (0, 1, 2, or 3 lines).
        \item \textbf{Arrows:} Point from the longer root to the shorter root.
    \end{itemize}

    \vfill
    % Bottone ipertestuale verso l'appendice per definizioni formali
    \hyperlink{backup_simpleroots}{\beamerbutton{Definition of Simple Roots}}
    \hyperlink{backup_dynkin}{\beamerbutton{Definition of Dynkin Diagrams}}
\end{frame}

\begin{frame}{The Main Result: Rank 2 Dynkin Diagrams}
    The geometric root systems directly translate into the fundamental Dynkin diagrams:

    \vspace{0.5cm}
    \begin{figure}
        \centering
        % INSERIRE QUI SCREENSHOT PAGINA 125 DEL PDF (Dynkin Diagrams for Rank 2)
        \textit{[Insert Image: Dynkin Diagrams $A_1 \oplus A_1$, $A_2$, $B_2$, $G_2$]}
    \end{figure}

    \vspace{0.5cm}
    From these simple building blocks, one can classify all semi-simple Lie algebras, setting the stage for constructing Grand Unified Theories (GUTs).
\end{frame}

